% ==========================================
% ACKNOWLEDGEMENTS
% ==========================================
\chapter*{Remerciements}

\addcontentsline{toc}{chapter}{Remerciements}

    Je tiens tout d'abord à remercier \textbf{Dafnis Krasniqi}, directeur de la formation, 
    professeur de deep learning et tuteur de ce mémoire, pour son encadrement, 
    ses retours précis et sa disponibilité. Mes remerciements s'adressent également 
    à \textbf{Marc Arthur Diaye}, ancien directeur de la formation.

    Je remercie l'ensemble de l'équipe pédagogique pour la qualité des enseignements 
    reçus : \textbf{Alexis Bogroff} (Python, machine learning), \textbf{Fatimetou Zeine} (SQL), 
    \textbf{Baptiste Gautier} (statistiques), \textbf{Ibrahima Diattara} (calcul distribué, 
    Spark, datastream), \textbf{Ansoumana Cisse} (Neo4j) et \textbf{Kamila Kare} (MLOps, IA générative).

    Une mention toute particulière est dédiée à ma sœur, \textbf{Lina}, également professeure 
    de machine learning au sein de ce cursus. Je la remercie de m'avoir recommandé 
    cette formation. Ses bons conseils et nos nombreuses discussions autour de la data 
    science et de l'intelligence artificielle ont grandement nourri ma réflexion.

    Enfin, je remercie ma famille pour son soutien constant tout au long de ce parcours.
    Je suis particulièrement reconnaissant envers ma mère, \textbf{Soumia}, pour son accompagnement et ses 
    encouragements. Une pensée toute spéciale pour mes deux petits garçons, \textbf{Ali} et \textbf{Adib}, 
    qui ont parfois dû patienter sagement.

\vspace{2cm} % Adds some vertical space if needed before the next section

% ==========================================
% GLOSSARY
% ==========================================
\chapter*{Glossaire}
\addcontentsline{toc}{chapter}{Glossaire}

\begin{acronym}[Prompt Engineering] 
    
    \acro{API}{Application Programming Interface\acroextra{\\
    Interface de programmation permettant à deux systèmes informatiques de communiquer 
    entre eux. Dans ce mémoire, désigne les interfaces d'accès aux modèles OpenAI.}}

    \acro{AST}{Abstract Syntax Tree\acroextra{\\
    Arbre syntaxique abstrait. Structure de données représentant le code source sous 
    forme d'arbre. Utilisé dans ce mémoire (via la bibliothèque sqlparse) pour analyser 
    la requête SQL générée et garantir un mode lecture seule (blocage des requêtes autres 
    que SELECT).}}

    \acro{Benchmark}{Benchmark\acroextra{\\
    Protocole d'évaluation standardisé permettant de comparer les performances de 
    différents systèmes sur un ensemble de tâches identiques et reproductibles.}}

    \acro{BIRD}{Big Bench for Large-Scale Database Grounded Text-to-SQL Evaluation\acroextra{\\
    Jeu de données d'évaluation Text-to-SQL introduisant des bases de données massives 
    avec des données bruitées pour simuler les conditions réelles d'entreprise.}}

    \acro{Chain-of-Thought}{Chain-of-Thought (CoT)\acroextra{\\
    Technique de prompt engineering forçant le modèle à détailler son raisonnement étape 
    par étape avant de produire une réponse finale. Traduit dans ce mémoire par "raisonnement 
    par étapes".}}

    \acro{CRUD}{Create, Read, Update, Delete\acroextra{\\
    Les quatre opérations fondamentales d'une base de données relationnelle : création, 
    lecture, mise à jour et suppression de données.}}

    \acro{DAG}{Directed Acyclic Graph\acroextra{\\
    Graphe orienté sans cycle. Utilisé dans les frameworks d'orchestration classiques 
    (Apache Airflow). Insuffisant pour modéliser les agents LLM nécessitant des boucles 
    itératives.}}

    \acro{DAIL-SQL}{DAIL-SQL\acroextra{\\
    Méthode Text-to-SQL optimisant la sélection dynamique des exemples Few-Shot par 
    similarité sémantique pour guider le modèle de langage.}}

    \acro{DDL}{Data Definition Language\acroextra{\\
    Sous-ensemble du langage SQL dédié à la définition de la structure d'une base de 
    données (instructions CREATE TABLE, ALTER TABLE, etc.).}}

    \acro{DIN-SQL}{DIN-SQL\acroextra{\\
    Méthode Text-to-SQL décomposant la tâche en quatre modules distincts : Schema 
    Linking, décomposition, génération et auto-correction.}}

    \acro{Embeddings}{Embeddings\acroextra{\\
    Représentations vectorielles numériques de textes permettant de mesurer leur 
    similarité sémantique. Utilisés dans ce mémoire pour la sélection dynamique des 
    exemples Few-Shot.}}

    \acro{ERD}{Entity-Relationship Diagram\acroextra{\\
    Diagramme entité-association représentant graphiquement la structure d'une base 
    de données et les relations entre ses tables.}}

    \acro{Few-Shot Prompting}{Few-Shot Prompting\acroextra{\\
    Technique consistant à injecter quelques exemples de paires question/réponse dans 
    l'instruction donnée au modèle pour guider sa génération. Traduit dans ce mémoire 
    par "injection d'exemples en contexte".}}

    \acro{Fine-Tuning}{Fine-Tuning\acroextra{\\
    Réentraînement d'un modèle de langage pré-entraîné sur un jeu de données spécifique 
    pour spécialiser ses capacités sur un domaine ou une tâche précise.}}

    \acro{FK}{Foreign Key\acroextra{\\
    Clé étrangère. Colonne d'une table référençant la clé primaire d'une autre table pour 
    établir une relation entre elles.}}

    \acro{Guardrails}{Guardrails\acroextra{\\
    Mécanismes de sécurité implémentés dans un système d'IA pour contraindre ou filtrer
    les sorties du modèle et prévenir les comportements indésirables.}}

    \acro{Human-in-the-loop}{Human-in-the-loop\acroextra{\\
    Paradigme d'évaluation ou d'entraînement intégrant un expert humain dans la boucle de 
    validation pour superviser ou corriger les décisions automatisées.}}

    \acro{JSON}{JavaScript Object Notation\acroextra{\\
    Format de données textuelles léger et lisible, utilisé dans ce mémoire pour stocker 
    le benchmark et l'historique des sessions d'évaluation.}}

    \acro{LangGraph}{LangGraph\acroextra{\\
    Bibliothèque Python de la suite LangChain permettant de modéliser des agents LLM 
    sous forme de graphes d'états cycliques.}}

    \acro{LLM}{Large Language Model\acroextra{\\
    Grand modèle de langage. Modèle d'intelligence artificielle entraîné sur de vastes 
    corpus textuels, capable de comprendre et de générer du langage naturel.}}

    \acro{LLM-as-a-Judge}{LLM-as-a-Judge\acroextra{\\
    Protocole d'évaluation automatisée utilisant un LLM puissant comme évaluateur pour 
    noter la qualité sémantique des réponses générées par d'autres modèles.}}

    \acro{LLMOps}{Large Language Model Operations\acroextra{\\
    Déclinaison du MLOps spécifiquement adaptée au cycle de vie des grands modèles de 
    langage, incluant la gestion des prompts, la surveillance de la dérive (Prompt Drift) 
    et l'évaluation continue (LLM-as-a-judge).}}

    \acro{MLOps}{Machine Learning Operations\acroextra{\\
    Ensemble de pratiques d'ingénierie visant à fiabiliser, automatiser et surveiller 
    le déploiement des modèles de machine learning en production.}}

    \acro{NLP}{Natural Language Processing\acroextra{\\
    Traitement automatique du langage naturel. Domaine de l'intelligence artificielle 
    étudiant l'interaction entre les ordinateurs et le langage humain.}}

    \acro{Olist}{Olist\acroextra{\\
    Jeu de données public brésilien de commerce en ligne utilisé comme cas d'usage 
    principal et environnement d'évaluation dans ce mémoire.}}

    \acro{PK}{Primary Key\acroextra{\\
    Clé primaire. Colonne ou ensemble de colonnes identifiant de manière unique 
    chaque ligne d'une table de base de données.}}

    \acro{Prompt}{Prompt\acroextra{\\
    Instruction textuelle fournie à un modèle de langage pour guider sa génération. 
    Traduit dans ce mémoire par "invite".}}

    \acro{Prompt Drift}{Prompt Drift\acroextra{\\
    Dégradation progressive des performances d'un système LLM due aux mises à jour 
    silencieuses des modèles par leur fournisseur, sans modification du prompt.}}

    \acro{Prompt Engineering}{Prompt Engineering\acroextra{\\
    Discipline consistant à optimiser la formulation des instructions données à un 
    modèle de langage pour améliorer la qualité de ses réponses. Traduit dans ce 
    mémoire par "ingénierie des invites".}}

    \acro{RAG}{Retrieval-Augmented Generation\acroextra{\\
    Architecture combinant la recherche documentaire et la génération de texte pour 
    enrichir les réponses d'un LLM avec des informations externes.}}

    \acro{ReAct}{Reasoning and Acting\acroextra{\\
    Paradigme agentique forçant le modèle à alterner entre phases de raisonnement, 
    d'action externe et d'observation du résultat.}}

    \acro{RLHF}{Reinforcement Learning from Human Feedback\acroextra{\\
    Apprentissage par renforcement basé sur les retours humains. Méthode permettant 
    d'aligner le comportement d'un modèle d'IA sur les attentes humaines. Mentionnée 
    dans les préconisations pour améliorer l'agent grâce aux évaluations des utilisateurs.}}

    \acro{RLS}{Row-Level Security\acroextra{\\
    Mécanisme de sécurité filtrant automatiquement les lignes accessibles dans une 
    base de données selon le profil de l'utilisateur connecté.}}

    \acro{Schema Linking}{Schema Linking\acroextra{\\
    Étape critique du pipeline Text-to-SQL consistant à identifier, parmi toutes les 
    tables et colonnes du schéma, celles pertinentes pour la question posée.}}

    \acro{Self-Correction}{Self-Correction\acroextra{\\
    Mécanisme agentique permettant à l'agent de détecter ses propres erreurs d'exécution 
    et de régénérer une réponse corrigée de manière autonome. Traduit dans ce mémoire par 
    "auto-correction".}}

    \acro{Spider}{Spider\acroextra{\\
    Jeu de données d'évaluation Text-to-SQL de référence (Yu et al., 2018), comprenant 
    10 181 questions annotées sur 200 bases de données couvrant 138 domaines différents.}}

    \acro{SQL}{Structured Query Language\acroextra{\\
    Langage de requête structuré permettant d'interroger et de manipuler des bases de 
    données relationnelles.}}

    \acro{SQLite}{SQLite\acroextra{\\
    Système de gestion de base de données relationnelle léger, stockant l'intégralité 
    de la base dans un seul fichier local. Utilisé dans ce mémoire pour héberger la 
    base Olist.}}

    \acro{Streamlit}{Streamlit\acroextra{\\
    Bibliothèque Python permettant de créer rapidement des applications web interactives 
    de visualisation de données.}}

    \acro{Text-to-SQL}{Text-to-SQL\acroextra{\\
    Tâche de traduction automatique d'une question en langage naturel en une requête 
    SQL exécutable sur une base de données relationnelle.}}

    \acro{Zero-Shot}{Zero-Shot\acroextra{\\
    Utilisation d'un modèle de langage sans aucun exemple préalable dans l'instruction, 
    en s'appuyant uniquement sur ses connaissances générales acquises lors de 
    l'entraînement.}}

\end{acronym}